\documentclass{reportkit}

\setreportkitleftheader{REPORTKIT / ALGORITHM PRIMITIVE TEST}
\setreportkitfooter{Algorithm primitive acceptance test}
\setreportkitversion{v1.10.0}

\title{ReportKit -- algorithm/pseudocode primitive}
\author{Test build}
\date{15 September 2026}

\begin{document}
\maketitle

\begin{execsummary}
This is an acceptance test for \texttt{algorithmblock}, the language-neutral
algorithm/pseudocode primitive added by reportkit-algorithms.sty. It is
deliberately distinct from \texttt{codeblock} (executable, monospace source)
and from \texttt{reportflow}/\texttt{reportstate} (process relationships):
this primitive is for control-flow pseudocode set in proportional type.
\end{execsummary}

\section{Basic sequence}
A short, unconditional sequence of statements with input/output metadata.

\begin{algorithmblock}[label={alg:sequence}]{Running total}
  \AlgorithmInput{A list of numbers $x_1,\ldots,x_n$}
  \AlgorithmOutput{The running total}
  \State $total \gets 0$
  \State $total \gets total + x_1$
  \State $total \gets total + x_2$
  \State \Return $total$
\end{algorithmblock}

\section{Nested conditionals}
Two-pointer elimination, the motivating example for this primitive: nested
\texttt{\string\If}/\texttt{\string\Else} inside a \texttt{\string\While}
loop, with mathematics inside the statements.

\begin{algorithmblock}[label={alg:twopointer},caption={Two-pointer elimination finds the maximum container area in linear time.}]{Two-pointer elimination}
  \AlgorithmInput{Heights $h_0,\ldots,h_{n-1}$}
  \AlgorithmOutput{Maximum container area}
  \State $L \gets 0$
  \State $R \gets n-1$
  \State $best \gets 0$
  \While{$L < R$}
    \State $best \gets \max(best,(R-L)\min(h_L,h_R))$
    \If{$h_L \le h_R$}
      \State $L \gets L+1$
    \Else
      \State $R \gets R-1$
    \EndIf
  \EndWhile
  \State \Return $best$
\end{algorithmblock}

\section{Nested loops}
Nested \texttt{\string\For} inside a brute-force enumeration, for contrast
with the linear-time version above.

\begin{algorithmblock}{Brute-force pair enumeration}
  \AlgorithmInput{Array $h$ of length $n$}
  \AlgorithmOutput{Maximum area}
  \State $best \gets 0$
  \For{$i \gets 0$ to $n-2$}
    \For{$j \gets i+1$ to $n-1$}
      \State $a \gets (j-i)\min(h_i,h_j)$
      \State $best \gets \max(best,a)$
    \EndFor
  \EndFor
  \State \Return $best$
\end{algorithmblock}

\section{ElseIf chains and a worked binary search}
\texttt{\string\ElsIf} chains render with the same keyword treatment as
\texttt{\string\If}/\texttt{\string\Else}.

\begin{algorithmblock}[label={alg:binarysearch}]{Binary search}
  \AlgorithmInput{Sorted array $A$, target $x$}
  \AlgorithmOutput{Index of $x$, or $-1$}
  \State $L \gets 0$
  \State $R \gets |A|-1$
  \While{$L \le R$}
    \State $M \gets \lfloor (L+R)/2 \rfloor$
    \If{$A_M = x$}
      \State \Return $M$
    \ElsIf{$A_M < x$}
      \State $L \gets M+1$
    \Else
      \State $R \gets M-1$
    \EndIf
  \EndWhile
  \State \Return $-1$
\end{algorithmblock}

\section{Long algorithm spanning a page boundary}
A deliberately long, repetitive algorithm to exercise the breakable
tcolorbox across a page break -- this primitive is a non-floating reading
unit, so it must be able to break rather than drift away from this
paragraph or overflow the page.

\begin{algorithmblock}{Iterative relaxation sweep}
  \AlgorithmInput{Graph $G=(V,E)$ with edge weights $w$, source $s$}
  \AlgorithmOutput{Shortest distance from $s$ to every vertex}
  \State $dist[s] \gets 0$
  \For{each vertex $v \ne s$}
    \State $dist[v] \gets \infty$
  \EndFor
  \For{$i \gets 1$ to $|V|-1$}
    \For{each edge $(u,v)$ with weight $w(u,v)$ in $E$}
      \If{$dist[u] + w(u,v) < dist[v]$}
        \State $dist[v] \gets dist[u] + w(u,v)$
      \EndIf
    \EndFor
  \EndFor
  \For{each edge $(u,v)$ with weight $w(u,v)$ in $E$}
    \If{$dist[u] + w(u,v) < dist[v]$}
      \State \Return \textsc{negative-cycle-detected}
    \EndIf
  \EndFor
  \For{each vertex $v$ in $V$}
    \State \Call{Emit}{$v$, $dist[v]$}
  \EndFor
  \For{each vertex $v$ in $V$}
    \State \Call{Audit}{$v$, $dist[v]$}
  \EndFor
  \For{each vertex $v$ in $V$}
    \State \Call{Archive}{$v$, $dist[v]$}
  \EndFor
  \State \Return $dist$
\end{algorithmblock}

\section{Grayscale and semantic check}
\begin{principle}{Pseudocode is not source code}
\texttt{algorithmblock} and \texttt{codeblock} follow the same visual grammar
(Surface fill, hairline rule, square corners, compact spacing, small
sans-serif title) so the two read as one family, but they are not
interchangeable: \texttt{codeblock} is for executable, monospaced source;
\texttt{algorithmblock} is for language-neutral control flow set in
proportional type, with keywords carrying structure instead of syntax
highlighting.
\end{principle}

\section{Code and output blocks, for contrast}
\begin{codeblock}[two-pointer.py]{Python}
def max_area(heights):
    left, right, best = 0, len(heights) - 1, 0
    while left < right:
        best = max(best, (right - left) * min(heights[left], heights[right]))
        if heights[left] <= heights[right]:
            left += 1
        else:
            right -= 1
    return best
\end{codeblock}

\begin{outputblock}
87
\end{outputblock}

\end{document}
